\documentclass[acmsmall,nonacm,screen]{acmart}

\usepackage{mathpartir}
\usepackage{booktabs}

\citestyle{acmauthoryear}
\settopmatter{printacmref=false}

\AtEndPreamble{\theoremstyle{acmdefinition}\newtheorem{remark}[theorem]{Remark}}

\newcommand{\kw}[1]{\mathsf{#1}}
\newcommand{\nm}[1]{\mathit{#1}}
\newcommand{\lamx}[1]{\lambda #1.\,}
\newcommand{\rec}{\kw{rec}}
\newcommand{\suc}{\kw{S}}
\newcommand{\refl}{\kw{refl}}
\newcommand{\Nat}{\kw{Nat}}
\newcommand{\Eq}[3]{\kw{Eq}_{#1}(#2, #3)}
\newcommand{\conv}{\equiv}
\newcommand{\chk}[4]{\kw{check}(#1, #2, #3, #4)}
\newcommand{\yes}{\kw{yes}}
\newcommand{\no}{\kw{no}}
\newcommand{\out}{\kw{out}}
\newcommand{\cost}[3]{\kw{cost}(#1, #2, #3)}
\newcommand{\adv}[1]{\textup{\texttt{#1}}}

% A sentence that states a fact in words. The script has one check for every name.
\newcommand{\claim}[2]{#2}

% The rows of a table that the script writes.
\makeatletter
\newcommand{\rows}[1]{\@@input artifact/expected/#1.tex }
\makeatother

% written by pays.py; do not edit
\newcommand{\exAdversaryLargest}{10}
\newcommand{\exBetter}{31}
\newcommand{\exBinders}{eight}
\newcommand{\exCaseThree}{52}
\newcommand{\exConvertible}{1,420}
\newcommand{\exEquations}{eight}
\newcommand{\exLibraryBlind}{374,525}
\newcommand{\exLibraryPrice}{8,550}
\newcommand{\exLibraryRatio}{43}
\newcommand{\exLines}{900}
\newcommand{\exLiteralPairs}{300}
\newcommand{\exMutants}{five}
\newcommand{\exOrderAdvice}{48}
\newcommand{\exOrderBad}{13,537}
\newcommand{\exOrderGood}{48}
\newcommand{\exOverviewAdvised}{48}
\newcommand{\exOverviewBlind}{122,034}
\newcommand{\exPairs}{3,000}
\newcommand{\exPolicies}{six}
\newcommand{\exRandomPolicies}{300}
\newcommand{\exRatio}{38.1}
\newcommand{\exSearchRuns}{40}
\newcommand{\exSearchSteps}{1,099,566}
\newcommand{\exSearched}{600}
\newcommand{\exSeed}{2026}
\newcommand{\exTheorems}{seven}
\newcommand{\exUserEquations}{200}


\begin{document}

\title{The Prover Pays: Advice and a Price for Conversion in Proof Checking}

\author{Claude-Fable-5.2}
\authornote{An AI model. The proofs are on paper and have not been checked by a machine. A
script implements the checker and tests the theorems on random terms. At the time of writing
nobody has reviewed the paper.}

\begin{abstract}
A proof in type theory may be a single word, $\refl$, and checking it may take longer than
anyone is willing to wait, because the checker has to find out by computation that two terms
are equal. How it computes is decided by a strategy that is built into the checker: which of
two definitions to unfold first, and whether to compare the arguments of a function before
looking inside it. A good strategy is fast on the proofs that its authors had in mind.
When proofs are written by programs and submitted in bulk, the time to check one becomes part
of what the proof costs, and the submitter is the one who knows how the computation should
go. We let the submitter say it. A proof is a term together with \emph{advice}, a string of
answers to the questions that the conversion test asks, and a \emph{price}, a number of steps
that the checker will not exceed. We prove that advice cannot change the verdict, only the
cost; that the checker stops within the price; and that advice cannot be replaced by a
strategy of the kind that kernels have: for every strategy that decides by the ranks of the
two constants in question there are equations of size $n$ that it checks in more than $2^n$
steps and that two letters of advice check in a number of steps linear in $n$. A library of
\exTheorems{} small theorems costs \exLibraryPrice{} steps with advice and
\exLibraryBlind{} without. A script tests the claims on \exPairs{} random pairs of terms
against a separate evaluator and notices \exMutants{} deliberate mistakes.
\end{abstract}

\keywords{type theory, conversion, proof checking, proof certificates, definitional equality, cost}

\authorsaddresses{}

\maketitle

\section{Introduction}

In a type theory with computation, the terms $2 + 2$ and $4$ are the same by definition, and
the proof that they are equal is $\refl$. The same holds for $2^{1000} \bmod 7$ and its
value, and for the result of running a verified decision procedure on a large formula. This
is what makes such theories practical: computation is not written into proofs. It is also
what makes the time to check a proof unpredictable. The size of a proof says nothing about
it. The checker has to decide whether two terms are \emph{convertible}, and the decision is
a computation whose length depends on the terms and on how the checker goes about it.

How it goes about it is a matter of strategy. When the two terms are applications of
different defined constants, the checker must unfold one of them, and unfolding the wrong
one can turn a comparison of two short terms into a comparison of two enormous ones. When
they are applications of the same constant, it can compare the arguments first, which is
fast when the arguments are equal and wasted when the function ignores them. Kernels decide
such questions by rules of thumb. The kernel of Rocq, for one, gives every constant a level
and unfolds ``the one with lower level'' first, and users can set the levels by a
command~\cite{rocqstrategy}. The kernels of Lean give every definition a height, unfold the
higher of two, and compare arguments first when the two are the same~\cite{bailey2024}. The rules work well on the proofs that their authors have seen.
Users of proof assistants know the other case: a proof that should check at once and does not
return, repaired by marking a definition opaque or by restating a goal.

This was a nuisance when proofs were written by people who wait for their own checker. It is
a different matter when proofs are produced by programs and checked by somebody else. A
registry that accepts machine-written proofs in bulk has to run each one, and a proof whose
check does not return is an attack on the registry, whether or not it was meant as one. The
time to check a proof becomes part of what the proof costs, and the question is who pays and
who decides how the computation is done.

Our answer is that the prover does both. The prover knows why the two terms are equal: it
constructed them. So it can tell the checker how to see it, and it can say in advance how
long the check will take.

\paragraph{The proposal.} A proof is a triple of a term, \emph{advice}, and a \emph{price}.
The advice is a string over four letters. The conversion test of the checker is the usual
one, with one change: at each point where a strategy would decide, it reads the next letter
of the advice. The price is a number of steps. The checker counts its steps and gives up
when the count passes the price. Nothing else is new, and the checker is, if anything,
simpler than before, since it contains no strategy.

\paragraph{The results.} Advice is not trusted and need not be. Whatever the advice, the
checker never accepts terms that are not convertible, and never rejects terms that are, as
long as the price suffices (Theorems~\ref{thm:sound} and~\ref{thm:advice}): advice changes
the cost and nothing else. The checker stops within the price (Theorem~\ref{thm:price}), so
a registry knows the cost of a submission before it runs it and can refuse what it does not
want to pay for. And advice does something that no strategy of the kernel kind can do. For
every strategy that decides by the ranks of the two constants, as the levels above do, we
construct definitions and equations of size $n$ on which the strategy takes more than $2^n$
steps while two letters of advice suffice for a number of steps linear in $n$
(Theorem~\ref{thm:adversary}). The strategy cannot know which of the two constants is the
expensive one. The prover does.

Our main contributions are as follows:
\begin{itemize}
\item The conversion test with advice and a price, for a type theory with functions,
  numbers with recursion, definitions, and equations proved by computation
  (Section~\ref{sec:test}).
\item The theorems that advice is only advice and that the price is kept
  (Theorems~\ref{thm:sound}, \ref{thm:advice}, and~\ref{thm:price}).
\item The theorem that no strategy by ranks can replace advice
  (Theorem~\ref{thm:adversary}), with an adversary that is a program: given a strategy, it
  returns the definitions and the equations.
\item Proofs as triples, with a checker for the type theory whose every use of conversion
  draws on one string of advice and one price, so that the price of a library is the sum of
  the prices of its theorems and a theorem that is used is not paid for again
  (Section~\ref{sec:proofs}).
\item An implementation, tested on \exPairs{} random pairs of terms against an evaluator
  that shares nothing with it (Section~\ref{sec:eval}).
\end{itemize}

The parts of the proposal are not new one by one. Proofs that carry hints for their checker
go back to proof-carrying code, where an oracle resolves the choices of a proof
search~\cite{necula2001}, and proofs of unsatisfiability are enriched with hints so that a
verified checker need not search~\cite{cruzfilipe2017}. Conversion has been made
explicit in the terms of pure type systems~\cite{vandoorn2013}. What we add is the
combination for conversion in a kernel: advice at exactly the points where kernels have a
strategy, a price, the proof that the advice is harmless, and the proof that it is needed.

\section{Overview}
\label{sec:overview}

Take numbers with a successor $\suc$ and a recursor, and these definitions: $\nm{id}$ is the
identity, $\nm{double}\ n$ adds $n$ to itself by recursion, and $\nm{pow2}\ n$ doubles $1$
for $n$ times. Consider the equation
\[
  \nm{pow2}\ 10 \;=\; \nm{id}\ (\nm{pow2}\ 10) .
\]
It holds by computation, and its proof is $\refl$. To check the proof, the checker must
decide that the two sides are convertible. Both sides are applications of a constant, and
the constants differ. A checker that unfolds both computes $\nm{pow2}\ 10$, a tower of
$2^{10}$ successors, then finds on the right $\nm{pow2}\ 10$ again after unfolding
$\nm{id}$, computes the second tower, and compares the two.
\claim{overview}{That takes \exOverviewBlind{} steps of the machine of this paper.}

A checker that unfolds the right side first sees $\nm{pow2}\ 10$ against $\nm{pow2}\ 10$: the
same constant with arguments that can be compared without unfolding anything. \claim{advice-rs}{That
takes \exOverviewAdvised{} steps.} The two decisions, \emph{unfold on the right} and then
\emph{compare the arguments}, are the advice \adv{RS}. The prover that produced the equation
knows them; the checker cannot. \claim{walk}{Table~\ref{tab:walk} shows both runs: what the
test asks, what the advice answers, and what is unfolded.} With \adv{RS} two questions are
asked and one constant is unfolded. With \adv{L} the first answer is wrong, the advice is
used up, and everything is unfolded.

\begin{table}
\caption{The test on $\nm{pow2}\ 10 = \nm{id}\ (\nm{pow2}\ 10)$ under two strings of advice.
The second column is the number of the step at which the event happens, and for the last
line of each run the number of steps in all.}
\label{tab:walk}
\begin{tabular}{lrl}
\toprule
advice & step & event \\
\midrule
\rows{walk}
\bottomrule
\end{tabular}
\end{table}

One might hope that a better rule of thumb does as well. The rule ``unfold the constant that
was defined later'' gets this example right if $\nm{id}$ was defined after $\nm{pow2}$ and
wrong if it was defined before. The rule ``compare arguments first'' is right here and wrong
for $\nm{const0}\ (\nm{pow2}\ 10) = \nm{const0}\ (\nm{pow2}\ 11)$, where $\nm{const0}$ ignores
its argument: the arguments are different numbers, and the checker computes both of them
before it gives up and unfolds $\nm{const0}$. Section~\ref{sec:need} shows that this is not
bad luck with these rules. Every rule of this kind has its example.

Advice may be wrong, by error or by intent. With the advice \adv{L} the checker of the
equation above does the slow thing. It still reaches the right verdict, if the price allows,
and if the price does not allow, it stops and says so. A proof with bad advice is a proof
that is expensive, and the price that comes with it says how expensive.

\section{Terms and Conversion}
\label{sec:terms}

The terms are those of the lambda calculus with numbers, a recursor, and constants:
\[
  t, u \;::=\; x \mid \lamx{x} t \mid t\ u \mid c \mid 0 \mid \suc\ t \mid \rec\ t\ u\ n
\]
A \emph{library} is a list of definitions $c := t$ with closed bodies, in which a body uses
only earlier constants. The \emph{rank} of a constant is its position in the library.
\emph{Conversion} $t \conv u$ is the least congruence that contains
\[
\begin{array}{r@{\;\conv\;}l@{\qquad}r@{\;\conv\;}l}
  (\lamx{x} t)\ u & t[u/x] & c & t \quad \text{for a definition } c := t \\
  \rec\ z\ s\ 0 & z & \rec\ z\ s\ (\suc\ n) & s\ n\ (\rec\ z\ s\ n) \\
  \lamx{x} t\ x & t \quad (x \text{ not free in } t)
\end{array}
\]
The first four equations, read from left to right, are the rules of reduction. We consider
terms that have types in the system of Section~\ref{sec:proofs}, which is the system T of
G\"odel with definitions and equations. For such terms reduction is confluent and every
reduction sequence ends~\cite{tait1967}, and two terms of the same type are convertible
exactly if their normal forms are equal up to the last equation.

\section{The Test}
\label{sec:test}

\subsection{The Machine}

The checker reduces terms to \emph{weak head normal form} on a machine with environments, in
the style of the algorithm of \citet{coquand1996}. A value is a closure $\langle \rho, \lamx{x}
t\rangle$, a number $0$ or $\suc\ \theta$, or a \emph{neutral} value $h\ \vec{e}$: a head that
is a variable or a constant, with a spine of eliminations, each the application to an
argument or a recursion. Arguments are thunks $\theta$: a term with an environment that is
evaluated when it is needed and at most once. The machine does not unfold constants by
itself. A constant evaluates to a neutral value with the constant as its head, and a separate
operation $\kw{unfold}$ replaces the head by the value of its definition and applies the
spine. This is what leaves the decision about unfolding to the test.

Every entry into the evaluator, every application, every step of a recursion, every
unfolding, and every comparison of the test counts as one \emph{step}. A step takes constant
time on a machine with pointers, except that the lookup of a variable takes time
proportional to its index~\cite{debruijn1972}.

\subsection{The Test with Advice}

The test compares two values, and we write $v \approx w$ when it succeeds. Its structural
part is standard and is given in Figure~\ref{fig:test}. Two numbers are compared by their
constructors. If one of the values is a closure, both are applied to a fresh variable, which
gives the last equation of conversion. Two neutral values with the same variable as head are
compared spine by spine, and two spines are compared elimination by elimination. Values that
are the same object are convertible without a further look, which is what makes a shared
argument cheap. When no rule applies, the test fails.

\begin{figure}
\begin{mathpar}
\inferrule{v \text{ and } w \text{ are the same object}}{v \approx w}
\and
\inferrule{ }{0 \approx 0}
\and
\inferrule{\theta{\downarrow} \approx \theta'{\downarrow}}{\suc\ \theta \approx \suc\ \theta'}
\and
\inferrule{v \text{ or } w \text{ is a closure} \\ v\ x \approx w\ x \quad (x \text{ fresh})}{v \approx w}
\and
\inferrule{\vec{e} \approx \vec{e}'}{x\ \vec{e} \approx x\ \vec{e}'}
\and
\inferrule{\theta_i{\downarrow} \approx \theta_i'{\downarrow} \text{ for the thunks of } e_i \text{ and } e_i', \text{ of one kind, for all } i}
          {e_1 \cdots e_k \approx e_1' \cdots e_k'}
\and
\inferrule{\kw{unfold}(c\ \vec{e}) \approx w \\ w \text{ has no constant at its head}}{c\ \vec{e} \approx w}
\and
\inferrule{v \approx \kw{unfold}(d\ \vec{e}') \\ v \text{ has no constant at its head}}{v \approx d\ \vec{e}'}
\end{mathpar}
\caption{The rules of the test that involve no choice. $\theta{\downarrow}$ is the value of a
thunk, computed at most once.}
\Description{Eight inference rules: identical objects, zero, successor, closures applied to a
fresh variable, neutral values with a variable at the head, spines, and the two rules that
unfold a constant when only one side has one at its head.}
\label{fig:test}
\end{figure}

The test has a choice exactly when both values are neutral with constants as heads.
Figure~\ref{fig:choice} gives the four things that it can do, and the advice says which. If
only one value has a constant as its head, that constant is unfolded and nothing is asked,
by the last two rules of Figure~\ref{fig:test}.

\begin{figure}
\[
\begin{array}{c@{\quad}l@{\quad}l}
\text{letter} & \text{the test continues with} & \text{allowed when} \\[2pt]
\adv{L} & \kw{unfold}(c\ \vec{e}) \text{ against } d\ \vec{e}' & \text{always} \\
\adv{R} & c\ \vec{e} \text{ against } \kw{unfold}(d\ \vec{e}') & \text{always} \\
\adv{B} & \kw{unfold}(c\ \vec{e}) \text{ against } \kw{unfold}(d\ \vec{e}') & \text{always} \\
\adv{S} & \vec{e} \text{ against } \vec{e}', \text{ and if they differ, as for } \adv{B} & c = d,\ \vec{e} \text{ and } \vec{e}' \text{ of one length}
\end{array}
\]
\caption{The choice of the test for $c\ \vec{e}$ against $d\ \vec{e}'$, where $c$ and $d$ are
constants. A letter \adv{S} where it is not allowed is read as \adv{B}, and so is the end of
the advice.}
\Description{A table of the four letters of advice: L unfolds the left constant, R the right
one, B both, and S compares the spines first and falls back to B if they differ.}
\label{fig:choice}
\end{figure}

The letter \adv{S} deserves a remark. If the spines are convertible, so are the two values,
by congruence. If they are not, nothing follows, because a function may give equal results
for different arguments, and the test goes on by unfolding. A comparison of spines that
fails has cost steps and decided nothing.

\begin{definition}[Check]
\label{def:check}
For closed terms $t$ and $u$, advice $\alpha$, which is a string of the four letters, and a
price $p$, $\chk{t}{u}{\alpha}{p}$ evaluates both terms, runs the test
with the letters of $\alpha$ in order, and returns $\yes$ or $\no$ if the test finishes
within $p$ steps and $\out$ otherwise. The cost $\cost{t}{u}{\alpha}$ is the number of steps
of the run without a price.
\end{definition}

\begin{theorem}[Soundness]
\label{thm:sound}
Let $t$ and $u$ be terms of one type. If $\chk{t}{u}{\alpha}{p} = \yes$ then $t \conv u$. If
$\chk{t}{u}{\alpha}{p} = \no$ then $t \not\conv u$.
\end{theorem}

\begin{proof}
Every value $v$ stands for a term $\lceil v \rceil$, its \emph{readback}: a closure for its
abstraction with the environment substituted, a thunk for its term likewise, a neutral value
for its head applied to the readbacks of its spine. The machine keeps conversion in this
sense: if a term $t$ evaluates in an environment $\rho$ to $v$, then $\lceil v \rceil \conv
t\rho$, and $\lceil \kw{unfold}(v) \rceil \conv \lceil v \rceil$. Both hold because every
step of the evaluator, of $\kw{unfold}$, and of the recursion is a reduction or the
bookkeeping of an environment, and because a thunk is evaluated in the environment in which
it was made. We show that $v \approx w$ implies $\lceil v \rceil \conv \lceil w \rceil$, by
induction on the run. The test answers $\yes$ only where
congruence, the last equation, or the identity of two values justifies it. For \adv{S} this
is the remark above. It answers $\no$ only for two values in weak head normal form with no
constant at the head that differ in their constructor, or that have different variables as
heads, or the same variable with spines that are not convertible. By confluence, terms with
such forms are not convertible. A comparison of spines under \adv{S} that fails does not
answer $\no$; the test unfolds and goes on.
\end{proof}

\begin{theorem}[Advice is advice]
\label{thm:advice}
Let $t$ and $u$ be terms of one type.
\begin{enumerate}
\item For every advice $\alpha$ the cost $\cost{t}{u}{\alpha}$ is finite.
\item For all advice $\alpha$ and $\alpha'$ and all prices $p$ and $p'$: if neither
  $\chk{t}{u}{\alpha}{p}$ nor $\chk{t}{u}{\alpha'}{p'}$ is $\out$, they are equal.
\item $\chk{t}{u}{\alpha}{p} = \out$ exactly if $p < \cost{t}{u}{\alpha}$.
\end{enumerate}
\end{theorem}

\begin{proof}
(1) Give a call of the test on $v$ and $w$ the measure that consists of the length of the
longest reduction sequence of $\lceil v \rceil$ and of $\lceil w \rceil$ together, then the
size of the two normal forms, and then one bit that says whether the letter \adv{S} is still
to be tried. An unfolding, an application of a closure, and a step of the recursion are
reductions and lower the first component, which is finite because the terms have types. A
descent into constructors or spines keeps it and lowers the second. A comparison of spines
under \adv{S} that fails is followed by an unfolding. So the run ends whatever the letters
are. (2) By Theorem~\ref{thm:sound} both verdicts say whether $t \conv u$. (3) The run with a
price is the run without, cut off when the count passes $p$.
\end{proof}

\begin{theorem}[The price is kept]
\label{thm:price}
$\chk{t}{u}{\alpha}{p}$ makes at most $p + 1$ steps, and reads at most $p$ letters of
$\alpha$.
\end{theorem}

\begin{proof}
The count is tested at every step, and reading a letter is part of a step.
\end{proof}

Advice cannot make a computation shorter than its result.

\begin{proposition}[The written answer]
\label{prop:literal}
If $u$ is the numeral for $m$, written out with $m$ successors, and $t$ is any term with $t
\conv u$, then $\cost{t}{u}{\alpha} \ge m$ for every advice $\alpha$.
\end{proposition}

\begin{proof}
The value of $u$ is built by the evaluator one successor at a time, each in a step of its
own, and only when the test asks for it. The test can answer $\yes$ only after it has
compared every successor of the value of $u$ with one of the value of $t$, or has found the
two values to be the same object, which they are not, since the successors of $u$ are made
by its own evaluation.
\end{proof}

\claim{literal}{For \exLiteralPairs{} random terms against their written values the bound holds
under every policy and under the cheapest advice.} The proposition marks the line between the
two parts of a cost. An equation between a program and its written result has a price that
is at least the size of the result, and no advice changes that. An equation between two
terms that are related by their shape, such as a term and the same term under an identity,
has a price that depends on seeing the shape, and that is what advice is for.

\paragraph{Other choices.} Unfolding is the choice that this test has, and it is the one that
decides between a short check and a long one in the examples that we know. A kernel for a
richer theory has more. It chooses the order in which the parts of two spines or two types
are compared, which matters when one part fails early; whether to try a cached answer; and,
with records and proofs that are irrelevant, when to apply the rules that are specific to
them. Each of these is a place for a further letter, and the theorems of this section do not
depend on how many letters there are: a letter selects among ways to continue that are each
sound, and all ways end.

The three theorems divide the work. The verdict belongs to the terms alone. The cost belongs
to the terms and the advice. The checker is responsible for neither a good strategy nor a
bound on its own running time: it follows and it counts. For a pair of convertible terms,
the least cost over all advice is a property of the pair, and we call it the \emph{price of
the equation}. Finding advice that attains it is a search, and the search is the prover's.

\section{Why Advice Is Needed}
\label{sec:need}

A \emph{policy} answers the question of Figure~\ref{fig:choice} without advice: it is a
function from the two constants, and from whether the letter \adv{S} is allowed, to a letter.
A policy \emph{decides by ranks} if it depends on the constants only through their ranks in
the library. The levels of the kernel mentioned in the introduction are ranks that the user
may reorder; heights of definitions and the order of declaration are ranks. Six policies of
this kind are: always \adv{L}; always \adv{R}; always \adv{B}; \emph{arguments first}, which
is \adv{S} where allowed and \adv{B} otherwise; and \emph{later first} and \emph{earlier
first}, which are \adv{S} where allowed and otherwise unfold the constant of the higher or
the lower rank.

\begin{theorem}[No policy by ranks]
\label{thm:adversary}
For every policy $\pi$ that decides by ranks there are a library and closed terms $t_n$ and
$u_n$ of the type of numbers, of size linear in $n$ and convertible, such that the test under
$\pi$ takes more than $2^n$ steps on $t_n$ and $u_n$, and there is advice of at most two
letters under which it takes a number of steps linear in $n$.
\end{theorem}

\begin{proof}
Let $\nm{big}$ be a closed term that maps the numeral $n$ to the numeral $2^n$ and contains
no constant, and let $\nm{thin}$ be $\lamx{x} x$. We ask $\pi$ four questions, in this order,
and build the library after its answers. Let $e$ and $l$ be two constants, $e$ of lower rank
than $l$, and let $\nm{const0}$ be a constant with the body $\lamx{x} 0$.

\emph{Case 1: $\pi$ answers \adv{S} for $\nm{const0}$ against itself.} Let $e := \nm{big}$,
$t_n = \nm{const0}\ (e\ n)$, and $u_n = \nm{const0}\ (e\ (n + 1))$. Under $\pi$ the test
compares $e\ n$ with $e\ (n+1)$. Whatever $\pi$ answers there, the arguments $n$ and $n + 1$
differ, both sides are unfolded, and the numerals $2^n$ and $2^{n+1}$ are compared, which
takes more than $2^n$ steps before they are found different. The advice \adv{B} unfolds
$\nm{const0}$ on both sides at once.

\emph{Case 2: otherwise, and $\pi$ does not answer \adv{S} for $e$ against itself.} Let
$e := \nm{big}$ and $t_n = u_n = e\ n$, written twice. Under $\pi$ one side is unfolded, the
other must follow, and two numerals $2^n$ are compared. The advice \adv{S} compares $n$ with
$n$.

\emph{Case 3: otherwise, and $\pi$ answers \adv{L} or \adv{B} for $e$ against $l$.} Let
$e := \nm{big}$, $l := \nm{thin}$, $t_n = e\ n$, and $u_n = l\ (e\ n)$. Under $\pi$ the left
side is unfolded to the numeral $2^n$. The right side then has a constant at its head and the
left side has not, so it is unfolded, twice, to the same numeral, and the numerals are
compared. The advice \adv{RS} unfolds $l$ and compares $n$ with $n$.

\emph{Case 4: otherwise.} Then $\pi$ answers \adv{R} for $e$ against $l$. Let
$e := \nm{thin}$, $l := \nm{big}$, $t_n = e\ (l\ n)$, and $u_n = l\ n$. Under $\pi$ the right
side is unfolded to the numeral, the left side must follow, and the numerals are compared.
The advice \adv{LS} unfolds $e$ and compares $n$ with $n$.

In each case the body of $\nm{big}$ has no constant, so that the test meets no further
choice once $\nm{big}$ is unfolded, and the comparison of two numerals takes at least as many
steps as the smaller one has successors. Under the advice named, the test compares the
numerals $n$ or constructors of constant size, after a constant number of unfoldings.
\end{proof}

The proof uses that the policy cannot tell which of two ranks holds the expensive
definition. A policy that reads the bodies of the definitions escapes this proof and not the
problem: whether unfolding a constant is expensive is a question about the value of a
program. A checker could also try the letters of Figure~\ref{fig:choice} side by side and
stop when one branch finishes. That costs a factor of at most four for every question on the
cheapest path, which is exponential in the length of the advice that it replaces, and it is
the search that the prover does once so that no checker has to do it again.

\begin{table}
\caption{The equations of Theorem~\ref{thm:adversary} for six policies: the case of the proof,
the steps of the test under the policy for four sizes, and the steps under the cheapest
advice for the largest size.}
\label{tab:adversary}
\begin{tabular}{llrrrrr}
\toprule
policy & case & $n = 4$ & $n = 6$ & $n = 8$ & $n = \exAdversaryLargest$ & with advice \\
\midrule
\rows{adversary}
\bottomrule
\end{tabular}
\end{table}

\begin{table}
\caption{\exEquations{} equations: the steps of the test under the six policies, and the least
number of steps that we found with advice of at most five letters or with one of the
policies.}
\label{tab:equations}
\small
\setlength{\tabcolsep}{3pt}
\begin{tabular}{lrrrrrrr}
\toprule
equation & left & right & both & arguments & later & earlier & least \\
\midrule
\rows{equations}
\bottomrule
\end{tabular}
\end{table}

\paragraph{Without an adversary.} The equations of the proof are built against a policy. The
equations of Table~\ref{tab:equations} are not. They are what a user writes: a term against
itself under an identity, a unit law, a doubling. \claim{equations}{For five of the
\exEquations{}, the worst policy takes more than a hundred times the steps of the best
advice, at least three different policies are the best one for some line, and for one line
advice makes no difference at all.} The first two lines are the same equation with two
identities that differ only in where they stand in the library, and the policies that
decide by rank exchange their columns. The third line punishes the policies that compare
arguments first, by seven million steps to sixteen. The last two lines have a cost that no
advice removes: the two sides have to be computed, because the equation is about their
values. A price has both parts in it, the cost of the computation that the equation needs
and the cost of the way in which the checker goes about it, and advice removes the second.

\begin{table}
\caption{\exUserEquations{} random equations of seven kinds that users write, with sizes
between five and eight: for how many a policy is within twice the least number of steps that
we found, for how many it takes more than a hundred times that number, and its largest
ratio.}
\label{tab:users}
\begin{tabular}{lrrr}
\toprule
policy & within twice & over a hundred times & largest ratio \\
\midrule
\rows{users}
\bottomrule
\end{tabular}
\end{table}

Table~\ref{tab:users} repeats the comparison on \exUserEquations{} equations drawn at random
from seven such kinds: a term against itself under one or two identities, the unit law, a
doubling written in two ways, a power against the doubling of the power below, a function
that ignores its argument, and a function applied to a wrapped argument. \claim{users}{No
policy is within twice the least cost for all of them, and every policy is more than a
hundred times above it for some.} The two policies that decide by rank and compare arguments
first are the best, as one expects from the kernels that use them, and they are near the
least cost for about half of the equations.

\claim{adversary}{Table~\ref{tab:adversary} shows the equations that the adversary of the
proof, run as a program, builds for the six policies; the same program builds equations with
the same growth for each of \exRandomPolicies{} random policies that decide by ranks, of which
\exCaseThree{} fall under the third or fourth case of the proof.} The steps under a policy
grow by a factor of about four when $n$ grows by two, as the numerals do. The steps under
advice grow by six.

\section{Proofs with a Price}
\label{sec:proofs}

Conversion is used by a type checker. We give the smallest type theory in which proofs by
computation exist: types are $\Nat$, dependent function types, and equations between terms
of a type,
\[
  A, B \;::=\; \Nat \mid \Pi x{:}A.\, B \mid \Eq{A}{t}{u}
\]
and terms have the additional form $\refl$. There is no universe and no induction principle,
so that an equation is proved by $\refl$ or by applying a function that returns such a proof.
Figure~\ref{fig:typing} gives the rules. They are those of a logical framework~\cite{harper1993}
with numbers; a library now gives every constant a type, and the body of a definition is
checked against it.

\begin{figure}
\begin{mathpar}
\inferrule{x : A \in \Gamma}{\Gamma \vdash x : A}
\and
\inferrule{c : A := t \text{ in the library}}{\Gamma \vdash c : A}
\and
\inferrule{\Gamma, x : A \vdash t : B}{\Gamma \vdash \lamx{x} t : \Pi x{:}A.\, B}
\and
\inferrule{\Gamma \vdash t : \Pi x{:}A.\, B \\ \Gamma \vdash u : A}{\Gamma \vdash t\ u : B[u/x]}
\and
\inferrule{ }{\Gamma \vdash 0 : \Nat}
\and
\inferrule{\Gamma \vdash t : \Nat}{\Gamma \vdash \suc\ t : \Nat}
\and
\inferrule{\Gamma \vdash z : \Nat \\ \Gamma \vdash s : \Nat \to \Nat \to \Nat \\ \Gamma \vdash n : \Nat}
          {\Gamma \vdash \rec\ z\ s\ n : \Nat}
\and
\inferrule{\Gamma \vdash t : A \\ \Gamma \vdash u : A \\ t \conv u}{\Gamma \vdash \refl : \Eq{A}{t}{u}}
\and
\inferrule{\Gamma \vdash t : A \\ A \conv B}{\Gamma \vdash t : B}
\end{mathpar}
\caption{Typing. The two rules of the last line use conversion.}
\Description{Nine typing rules: variables, constants, abstraction, application, zero,
successor, recursion at the type of numbers, reflexivity for convertible terms, and
conversion of types.}
\label{fig:typing}
\end{figure}

\begin{definition}[Proof]
\label{def:proof}
A \emph{proof} of a type $A$ is a triple of a term $t$, advice $\alpha$, and a price $p$. The
checker is bidirectional: it checks abstractions and $\refl$ against the given type and
infers the types of the other terms. Where the rules of Figure~\ref{fig:typing} use
conversion, it runs the test of Section~\ref{sec:test}. All runs of the test in the check of
one proof read from the one string $\alpha$, in order, and count against the one price $p$,
as does every other step of the checker. The proof is \emph{accepted} if the check ends
without a refusal within $p$ steps.
\end{definition}

\begin{theorem}[Proofs]
\label{thm:proofs}
If the proof $(t, \alpha, p)$ of $A$ is accepted, then $\vdash t : A$. Whether $t$ is
accepted at some price does not depend on $\alpha$. The check takes at most $p + 1$ steps.
A library whose definitions are checked one after the other is checked in at most the sum
of their prices and one step for each, and the check of a definition does not repeat the
check of an earlier one.
\end{theorem}

\begin{proof}
The checker follows the rules of Figure~\ref{fig:typing}, and by Theorem~\ref{thm:sound} it
uses the two rules with conversion only for convertible terms. A use of the test that
answers $\no$ refuses the proof, and by Theorem~\ref{thm:advice} that answer does not depend
on the advice. The steps are counted as in Theorem~\ref{thm:price}. The rule for a constant
reads the type of the constant from the library and does not look at its body.
\end{proof}

\begin{table}
\caption{A library of \exTheorems{} theorems: the advice, the price with it, and the steps of
the check when the advice is empty, which unfolds both sides at every choice.}
\label{tab:theorems}
\begin{tabular}{llrr}
\toprule
theorem & advice & price & without advice \\
\midrule
\rows{theorems}
\bottomrule
\end{tabular}
\end{table}

Table~\ref{tab:theorems} shows a small library. $\nm{four}$ is $2 + 2 = 4$ and
$\nm{sixteen}$ is $2^4 = 4 \cdot 4$; they are cheap either way. $\nm{wrapped}$ is the
equation of Section~\ref{sec:overview}. $\nm{ignored}$ is $\nm{const0}\ (\nm{pow2}\ 10) =
\nm{const0}\ (\nm{pow2}\ 11)$, for which unfolding both sides happens to be right.
$\nm{parity}$ says that $2^6$ is even, by running a function that decides evenness; it is a
proof by reflection in miniature, its cost is the computation, and no advice is needed or
useful. $\nm{unit}$ is the function that proves $0 + n = n$ for every $n$, by $\refl$ under a
binder. $\nm{use}$ applies it to $\nm{pow2}\ 10$: the type of the application and the type
that is asked for are the same equation, written twice, and the checker has to see that.
With the advice \adv{SSS} it compares the two types part by part without unfolding
$\nm{pow2}$. Without advice it unfolds, and compares numerals of $2^{10}$ successors, twice.
\claim{theorems}{The library costs \exLibraryPrice{} steps with its advice and
\exLibraryBlind{} without, \exLibraryRatio{} times as much, each proof is accepted at exactly
its price and not below it, and a proof of a false equation is refused whatever its advice.}

\paragraph{Who pays.} Consider a registry that receives proofs from provers that it does not
know.

\begin{proposition}[Budget]
\label{prop:budget}
A registry that checks proofs $(t_i, \alpha_i, p_i)$ in turn takes at most $\sum_i (p_i + 1)$
steps, whatever the proofs are. If it refuses every proof with $p_i > P$ before it runs it,
no submission costs it more than $P + 1$ steps, and the proofs that it refuses for this
reason are exactly those whose provers did not claim that $P$ steps suffice.
\end{proposition}

\begin{proof}
By Theorem~\ref{thm:proofs}, and because the price is read before the check starts.
\end{proof}

Without prices the registry has no such bound: a proof of one word may run for longer than
the registry exists, and a time limit that the registry sets by itself refuses proofs that
are correct and merely checked in the wrong way. A price is a promise by the prover and a
limit for the checker. A
registry can publish what it charges for a step, or the largest price that it accepts, and
refuse a submission after reading one number. A prover that understates the price has its
proof refused after the stated number of steps, and one that overstates it pays for steps
that are not taken. Advice that is wrong makes the proof more expensive than it had to be
and nothing worse. The reader of a library sees, next to each theorem, what it costs to
believe it, which is a fact about a proof that its size does not show~\cite{tpp1}.

\paragraph{Prices are stable.} A proof with advice is checked in the same way wherever it
stands. A proof under a policy is not.

\begin{proposition}[Stability]
\label{prop:stable}
The cost of checking a proof with its advice depends on the proof and on the definitions that
its terms reach, and not on the ranks of these definitions or on the rest of the library. For
each of the policies \emph{later first} and \emph{earlier first} there are two libraries that
differ only in the order of two definitions, and an equation of size $n$, such that the test
under the policy takes more than $2^n$ steps in one library and a number of steps linear in
$n$ in the other.
\end{proposition}

\begin{proof}
The test with advice consults nothing but the two values and the next letter. For the
second claim take the third and fourth cases of the proof of Theorem~\ref{thm:adversary}:
the equation of a term with a wrapper around itself, where the wrapper is defined after the
expensive function in one library and before it in the other.
\end{proof}

\claim{order}{For $n = 10$ the test under \emph{later first} takes \exOrderGood{} steps when
the wrapper is defined after the function that it wraps and \exOrderBad{} when it is defined
before; under \emph{earlier first} the numbers are exchanged; under the advice \adv{RS} it
takes \exOrderAdvice{} steps in both libraries.} A library that is reorganized, or a
definition that is given another level for the sake of one proof, can make other proofs slow
that nobody touched. Users of proof assistants know the effect. With advice, what a proof
costs is settled when the proof is made.

\paragraph{Finding advice.} The prover finds advice by running the test itself, with any
strategy it likes or with a search. The search over all advice of a given length is
exponential in the length, and the lengths are short: the advice of
Table~\ref{tab:theorems} has at most three letters. \claim{search}{For the equation of
Section~\ref{sec:overview}, a search that tries the letters in a fixed order and cuts every
run off at the best cost found so far runs the test \exSearchRuns{} times, for
\exSearchSteps{} steps in all, and returns \adv{RS}.} That is the cost of nine checks without
advice, paid once by the prover, against a saving at every later check. A prover that has constructed two terms
to be convertible, by a rewriting step or by a decision procedure, usually knows the
unfolding that makes it so.

\section{Evaluation}
\label{sec:eval}

The artifact is one Python script of about \exLines{} lines with no dependencies. It
implements the machine, the test with advice and a price, six policies, the search for the
cheapest advice, the adversary of Theorem~\ref{thm:adversary}, and the checker of
Section~\ref{sec:proofs}. For reference it has a second evaluator, which maps a term to a
Python number or function and shares nothing with the machine. The command \texttt{make -C
artifact check} runs the script and compares its output with the files from which this
paper takes its numbers and the rows of its tables.

\paragraph{Verdicts.} The script draws \exPairs{} pairs of closed terms of the type of
numbers from a generator with seed \exSeed{}, over a library with addition, doubling,
multiplication, powers of two, a predecessor, two identities, and a constant function;
\exConvertible{} of the pairs are convertible. \claim{verdicts}{For every pair the verdict of
the test is that of the reference, under each of the \exPolicies{} policies and under three
random strings of advice; with a price below the cost the answer is $\out$, after exactly
one step more than the price, and with a price equal to the cost it is the verdict}
(Theorems~\ref{thm:sound}, \ref{thm:advice}, and~\ref{thm:price}). \claim{binders}{For
\exBinders{} pairs of functions, among them pairs that are convertible only by the last
equation of conversion and a pair that is equal only by induction, the test gives the
expected verdict in both directions under every policy.} \claim{sharing}{A value that is
needed a second time costs no further step.}

\paragraph{Policies and advice.} For \exSearched{} of the random pairs the script searches
all advice of at most six letters for the cheapest. On random terms, which are small, the
policies do well: the cheapest advice beats every policy for \exBetter{} of the pairs, and the
worst ratio between a policy and the cheapest advice is \exRatio{}. The equations of
Table~\ref{tab:adversary} are not found by chance; they are built.

\paragraph{Deliberate mistakes.} The script plants \exMutants{} mistakes in the checker and
runs its tests against each. A checker that answers $\yes$ when the price runs out is
unsound. One that answers $\no$ when a comparison of spines under \adv{S} fails, or when only
the right value has a constant at its head, rejects convertible terms. One that does not
apply the last equation of conversion misses functions that are equal. One that evaluates a
thunk every time it is needed gives the right verdicts at a higher cost.
\claim{all-noticed}{Each of the five is noticed.}

\paragraph{What the tests do not show.} The reference decides conversion only for closed
terms of the type of numbers, where it is equality of values. For open terms and functions
the tests rest on a few chosen pairs and on the agreement between policies. The steps are
steps of this machine, which is an interpreter in Python; the numbers compare ways to check
on one machine and say nothing about the speed of a compiled kernel.

\section{Related Work}
\label{sec:related}

\paragraph{Proofs that help their checker.} Proof-carrying code ships a proof with a program
so that the receiver checks instead of trusting~\cite{necula1997}. \citet{necula2001} replace
the proof by an oracle, a string of bits that resolves the choices of a nondeterministic
checker, which makes proofs far smaller. Our advice is an oracle in this sense for one part
of a checker, conversion, and it is added to a proof and not put in its place.
In propositional satisfiability, proofs in the format of \citet{wetzler2014} are checked by
a search, and \citet{cruzfilipe2017} add hints that name the clauses each step uses, so that
a checker which is itself verified follows the proof without searching. The division of labor is
ours: the untrusted side searches once, the trusted side follows. An earlier paper of this
series studies programs that generate and check under a budget~\cite{tpp3}; a price is such
a budget, set by the party that generated.

\paragraph{Conversion in kernels.} Conversion by evaluation to weak head normal form with a
structural comparison is the algorithm of \citet{coquand1996}, and \citet{abel2017} prove a
version of it correct for a type theory with a universe. Kernels speed up the evaluation by
compilation~\cite{gregoire2002,boespflug2011}, which makes each step cheaper and leaves the
choice of steps where it was. The choice is made by levels, heights, and flags that users
can set~\cite{rocqstrategy,bailey2024}, which is advice of a kind, given once for a constant and not for
an equation; Theorem~\ref{thm:adversary} is about what such advice cannot express. In the
other direction, conversion can be removed from the checker altogether by making every
conversion step a part of the term: \citet{vandoorn2013} give pure type systems with
explicit convertibility proofs and show them equivalent to the usual
ones~\cite{adams2006,siles2012}, and frameworks based on rewriting make the rules of
computation a parameter~\cite{cousineau2007}. A term with all its conversion steps is as
large as the computation is long. Advice with a price keeps the term and the computation
apart, records only the decisions, and states the length.

\paragraph{Cost.} That the steps of an abstract machine are a sound measure of time is the
subject of a line of work on cost semantics~\cite{blelloch1996,accattoli2014}; our count is
of this kind, with the caveat about variable lookup. \citet{niu2022} make cost a part of a
type theory, for the programs that are its terms. We are concerned with the cost of checking,
which is a cost of the judgment and not of the term.

\paragraph{This series.} The first paper of the series asks what a person must read to
believe a machine-checked library~\cite{tpp1}. The price of a proof is a further item of
that kind: what a machine must spend. The practice of binding numbers and marked sentences
to the script follows~\cite{tpp10}.

\section{Limits and Conclusion}
\label{sec:limits}

\paragraph{Limits.} The type theory is small. It has no universes, no inductive families,
and no dependent elimination, and its only proofs of equations are by computation. The test
is the same for richer theories, with more constructors to compare, but we have neither
stated nor proved the theorems for them, and Theorem~\ref{thm:advice} needs that every
reduction sequence ends. The advice is tied to one algorithm: its letters answer the
questions that this test asks, in the order in which it asks them, and a different test
would need different advice. Proofs that carry advice are therefore less portable than
terms, although a checker can ignore advice and use a policy. The price is a number of steps
of one machine, and two implementations will agree on it only if they agree on what a step
is. We have not studied how hard it is to find the cheapest advice; the example of
Section~\ref{sec:need} suggests that it is at least as hard as knowing the value of a
program. Theorem~\ref{thm:adversary} is about policies that decide by ranks, and a policy
that inspects definitions is outside it. And the comparison of the two sides of an equation
that is false costs what it costs: advice shortens the way to $\yes$, and we have no theorem
about the way to $\no$.

\paragraph{Conclusion.} A kernel that decides how to compute has to be right for every proof
that it will ever see. A kernel that is told how to compute has to be right for none: it
follows, counts, and stops. The knowledge of how an equation holds is with whoever made the
equation, and a proof format should have a place for it, next to a number that says what the
proof costs to check.

\bibliographystyle{ACM-Reference-Format}
\bibliography{refs}

\end{document}
